\section{From vanishing slack to exponential products}\label{sec:rigidity}

Write $\Dgap(f)=h(f)-m-\tfrac12\log\det\Cov(f)$ for the
nonnegative entropy gap.  Theorem~\ref{mat:entropy-smooth} controls the nonlinear part of the normalized
matrix field and the sum of squares of its linear coefficients.  We now
show that this control determines every equality case.  The structural
step is local: the first two nonconstant terms of $q(A(x))$ force the
coefficients of a linear symmetric matrix field $A$ to commute whenever
$q(A)$ is a Jacobian.

\subsection{Compatibility of a nonlinear Jacobian}

\begin{lemma}[Rigidity of a linear matrix field]\label{lem:linear-jacobian}
Let $m\ge1$, and let $X_1,\ldots,X_m$ be real symmetric $m\times m$ matrices satisfying
\[
  \sum_{r=1}^m X_r^2=I,
  \qquad A(x)=\sum_{r=1}^m x_rX_r.
\]
Suppose $F\in W^{1,2}_{\mathrm{loc}}(\R^m;\R^m)$ has weak Jacobian
$DF=q(A)$ almost everywhere.  Then there are an orthogonal matrix $U$,
signs $\sigma_1,\ldots,\sigma_m\in\{-1,1\}$, and $w\in\R^m$ such that
\begin{equation}\label{eq:rigid-map}
  U^{\mathsf T}F(Uy)
  =\bigl(\sigma_i t(\sigma_i y_i)\bigr)_{i=1}^m+w
  \qquad\text{for almost every }y\in\R^m.
\end{equation}
In particular, $F_\#\gamma_m$ is an affine image of the product of $m$
mean-one exponential laws, and its covariance matrix is $I$.
\end{lemma}

\begin{proof}
Write $e_1,\ldots,e_m$ for the standard basis and $J=q(A)$.
Commutation of distributional derivatives of $F$ gives
\begin{equation}\label{eq:jacobian-compatibility}
  (\partial_rJ)e_s=(\partial_sJ)e_r
  \qquad(1\le r,s\le m).
\end{equation}
The scalar function $q$ is real analytic near zero.  Since $A(0)=0$, its
convergent power series, interpreted by matrix multiplication, gives
\[
  J(x)=q(0)I+q'(0)A(x)+\tfrac12q''(0)A(x)^2+O(|x|^3)
\]
near zero; the differentiated remainder is $O(|x|^2)$.
The identity $q'(a)=q(a)(q(a)-a)$ yields
\[
  q(0)=\sqrt{\frac2\pi},\qquad
  q'(0)=\frac2\pi>0,\qquad
  q''(0)=\sqrt{\frac2\pi}\left(\frac4\pi-1\right)>0.
\]
Thus \eqref{eq:jacobian-compatibility}, which is a pointwise identity
between analytic functions near zero, has constant term
\begin{equation}\label{eq:coefficient-symmetry}
  X_re_s=X_se_r.
\end{equation}
Equivalently, $A(x)z=A(z)x$ for all $x,z\in\R^m$; in particular,
$A(x)e_s=X_sx$.  Its degree-one term is
\[
  (X_rA(x)+A(x)X_r)e_s
  =(X_sA(x)+A(x)X_s)e_r.
\]
The terms with leftmost factor $A(x)$ cancel by
\eqref{eq:coefficient-symmetry}.  The remaining equality is
$(X_rX_s-X_sX_r)x=0$ near zero, and hence
\begin{equation}\label{eq:coefficients-commute}
  X_rX_s=X_sX_r\qquad(1\le r,s\le m).
\end{equation}

Commuting real symmetric matrices admit a common orthonormal eigenbasis.
Choose the orthogonal matrix $U$ with those basis vectors as columns,
and rotate both the input and the output.  Then
\[
  \widetilde A(y)=U^{\mathsf T}A(Uy)U
     =\sum_{i=1}^m y_i\widetilde X_i,
  \qquad
  \widetilde X_i=U^{\mathsf T}
      \left(\sum_{r=1}^m U_{ri}X_r\right)U
\]
has diagonal coefficients.  Orthogonality of $U$ gives
$\sum_i\widetilde X_i^2=I$, and the identity $A(x)z=A(z)x$ gives
$\widetilde X_i e_j=\widetilde X_j e_i$.  If $i\ne j$, these two vectors
are multiples of different basis vectors, so both are zero.
Consequently, only the $i$th diagonal entry of $\widetilde X_i$ can be
nonzero.  The sum-of-squares identity makes that entry a sign:
\[
  \widetilde X_i=\sigma_i e_ie_i^{\mathsf T},
  \qquad
  \widetilde A(y)=\operatorname{diag}(\sigma_1y_1,\ldots,\sigma_my_m).
\]
It follows that
\[
  D_y\bigl(U^{\mathsf T}F(Uy)\bigr)
   =\operatorname{diag}\bigl(q(\sigma_1y_1),\ldots,q(\sigma_my_m)\bigr).
\]
Since $t'=q$ and $\sigma_i^2=1$, the right side is also the Jacobian of
$y\mapsto(\sigma_i t(\sigma_i y_i))_{i=1}^m$.
Their difference has zero weak gradient on the connected set $\R^m$
and is therefore almost everywhere constant.  This proves
\eqref{eq:rigid-map}.

For independent standard Gaussian coordinates $G_i$, the variables
$t(\sigma_iG_i)$ are independent mean-one exponentials.  Their variances
are one, and the signs and the orthogonal matrix preserve covariance
$I$.  This proves the final assertion.
\end{proof}

\subsection{Compactness and the proof of the equality classification}

We first record the functional-calculus estimate that turns convergence
of the normalized fields into convergence of transport Jacobians.  The
Frobenius norm is essential to this elementary argument.

\begin{lemma}\label{lem:spectral-lipschitz}
For real symmetric matrices $M,N\in\R^{m\times m}$,
\begin{equation}\label{eq:spectral-lipschitz}
  |q(M)-q(N)|_{\mathrm F}\le |M-N|_{\mathrm F}.
\end{equation}
\end{lemma}

\begin{proof}
Let $(u_i)$ and $(v_k)$ be orthonormal eigenbases of $M$ and $N$, with
eigenvalues $\lambda_i$ and $\mu_k$.  Then
\[
  u_i^{\mathsf T}(q(M)-q(N))v_k
   =(q(\lambda_i)-q(\mu_k))\,u_i^{\mathsf T}v_k,
\]
whereas the corresponding entry of $M-N$ is
$(\lambda_i-\mu_k)u_i^{\mathsf T}v_k$.  The scalar bound $0<q'<1$
therefore bounds the absolute value of each former entry by the absolute
value of the latter.
Squaring and summing over $i,k$ proves \eqref{eq:spectral-lipschitz},
because the Frobenius norm is unchanged by orthogonal changes of the
row and column bases separately.
\end{proof}

\begin{proof}[Proof of the equality statement in Theorem~\ref{intro:entropy}]
Let $f$ be a log-concave probability density on $\R^m$ with
$\Dgap(f)=0$.  Choose the approximating densities $f_j$ from
Lemma~\ref{geo:approximation}.  Write $a_j,a$ for their means and the mean of
$f$, and $C_j,C$ for their covariance matrices.  That lemma gives
\begin{equation}\label{eq:equality-approximation}
  \norm{f_j-f}_{L^1(\R^m)}\longrightarrow0,\qquad
  a_j\longrightarrow a,\qquad C_j\longrightarrow C,\qquad
  \Dgap(f_j)\longrightarrow0.
\end{equation}
Here $C$ is positive definite, as recalled in the introduction, so the
covariance convergence also implies convergence of the log determinants.

For each $j$, apply Proposition~\ref{tr:normalization} to $f_j$.
Let $P_j$ be its positive definite normalizing matrix and let $F_j$
transport $\gamma_m$ to $(P_j)_\#(f_j(x)\,\dd x)$.  In the notation
of Theorem~\ref{mat:entropy-smooth}, put
\[
 \begin{gathered}
  J_j=DF_j,\qquad A_j=q^{-1}(J_j),\qquad
  X_{j,r}=\E[x_rA_j(x)],\\
  A_j(x)=\sum_{r=1}^m x_rX_{j,r}+Y_j(x),\qquad
  B_j=\sum_{r=1}^m X_{j,r}^2.
 \end{gathered}
\]
The entropy gap is affine invariant.  Theorem~\ref{mat:entropy-smooth} and
\eqref{eq:equality-approximation} therefore imply
\begin{equation}\label{eq:vanishing-slack}
  \norm{Y_j}_{L^2(\gamma_m)}\longrightarrow0,
  \qquad
  \sum_{i=1}^m \omega(\lambda_i(B_j))\longrightarrow0.
\end{equation}
The function $\omega$ is continuous and nonnegative on $[0,\infty)$,
vanishes only at $1$, and tends to infinity at infinity.
For every $\varepsilon>0$, it consequently has a positive infimum on
$\{\lambda\ge0:|\lambda-1|\ge\varepsilon\}$.  Thus every eigenvalue
of $B_j$ tends to $1$, and $B_j\to I$ in operator norm.

Since $\sum_r|X_{j,r}|_{\mathrm F}^2=\tr B_j$, the coefficients form
a bounded sequence in a finite-dimensional space.  Pass to a subsequence
such that $X_{j,r}\to X_r$ for every $r$.  Each $X_r$ is symmetric,
and Gaussian orthogonality together with \eqref{eq:vanishing-slack}
gives
\begin{equation}\label{eq:linear-field-limit}
  \sum_{r=1}^mX_r^2=I,\qquad
  A_j\longrightarrow A(x):=\sum_{r=1}^m x_rX_r
      \quad\text{in }L^2(\gamma_m).
\end{equation}
Lemma~\ref{lem:spectral-lipschitz} now gives
\begin{equation}\label{eq:jacobian-limit}
  J_j=q(A_j)\longrightarrow J:=q(A)
      \quad\text{in }L^2(\gamma_m).
\end{equation}

To obtain a limiting map, remove the irrelevant translations:
$\widehat F_j=F_j-\E F_j$.  Each $\widehat F_j$ lies in
$H^1(\gamma_m;\R^m)$: its law has finite second moments and, for this
fixed $j$, its Jacobian has the bounds \eqref{tr:Jacobian-bounds}
applied to the transformed target.  Gaussian Poincar\'e, applied
componentwise to the centered difference, gives
\begin{equation}\label{eq:map-cauchy}
  \norm{\widehat F_j-\widehat F_k}_{L^2(\gamma_m)}^2
    \le \norm{J_j-J_k}_{L^2(\gamma_m)}^2.
\end{equation}
Consequently $\widehat F_j$ converges strongly in Gaussian $H^1$ to a
centered map $\widehat F$ with weak Jacobian $D\widehat F=J$.
Indeed, \eqref{eq:map-cauchy} and \eqref{eq:jacobian-limit} give
convergence of both maps and first derivatives; passage to weak
derivatives is also valid on every bounded set, where the Gaussian
density has a positive lower bound.  In particular,
$\widehat F\in W^{1,2}_{\mathrm{loc}}(\R^m;\R^m)$.
Lemma~\ref{lem:linear-jacobian} applies to \eqref{eq:linear-field-limit}
and shows that $\widehat F_\#\gamma_m$ is an affine exponential product
with covariance $I$.

It remains to return from the normalized laws to $f$.
By Cauchy--Schwarz, the strong $L^2(\gamma_m)$ convergence of the
centered maps gives
\begin{equation}\label{eq:normalized-covariance-limit}
  D_j:=\Cov((\widehat F_j)_\#\gamma_m)
      =P_jC_jP_j^{\mathsf T}\longrightarrow I.
\end{equation}
We can now bound the normalizing matrices uniformly, without any
uniform Hessian bounds on the approximating potentials.
Choose positive constants $c,C_0,d,D$ such that, for all sufficiently
large $j$,
\[
  cI\le C_j\le C_0I,\qquad dI\le D_j\le DI.
\]
Congruence by $P_j$ then yields
\begin{equation}\label{eq:normalization-compactness}
  \frac d{C_0}I\le P_jP_j^{\mathsf T}\le\frac Dc I.
\end{equation}
Hence $P_j$ and $P_j^{-1}$ are bounded.  Pass to a further subsequence
with $P_j\to P$; the lower bound in
\eqref{eq:normalization-compactness} makes $P$ invertible.

The law of $\widehat F_j$ is the image of $f_j(x)\,\dd x$ under
$x\mapsto P_j(x-a_j)$.  For a bounded continuous function $\psi$,
\begin{align*}
 &\left|\int\psi(P_j(x-a_j))f_j(x)\,\dd x
       -\int\psi(P(x-a))f(x)\,\dd x\right|\\
 &\quad\le\norm{\psi}_{\infty}\norm{f_j-f}_{L^1}
     +\int\left|\psi(P_j(x-a_j))-\psi(P(x-a))\right|f(x)\,\dd x
     \longrightarrow0.
\end{align*}
The first term tends to zero by \eqref{eq:equality-approximation};
the second does so by dominated convergence.  Strong convergence of
$\widehat F_j$ also identifies the weak limit of these laws as
$\widehat F_\#\gamma_m$.  It follows that the image of $f(x)\,\dd x$
under $x\mapsto P(x-a)$ is the affine exponential product already
determined above.  Since $P$ is invertible, $f$ itself has the form
asserted in Theorem~\ref{intro:entropy}, up to equality almost everywhere.

Conversely, the density of a product of $m$ independent mean-one
exponentials is log-concave, has entropy $m$, and has covariance $I$.
Every invertible affine image remains log-concave and has the same
entropy gap by affine invariance.  Thus every density in the stated
class satisfies $\Dgap=0$.
\end{proof}
